\documentclass[11pt]{article}
\usepackage[margin=1in]{geometry}
\usepackage{amsmath,amssymb}
\usepackage{booktabs}
\usepackage{hyperref}

\title{nightsearch-SAST: Minimal Synthetic Dictionary Baseline for Cross-Attention Spot Annotation}
\author{nightsearch-SAST Project}
\date{\today}

\begin{document}
\maketitle

\section{Implemented Component}
This update adds a runnable \emph{synthetic dictionary baseline} to support end-to-end
experimentation for spot annotation using cross-attention. The implemented pipeline includes:
\begin{enumerate}
    \item typed config support for synthetic data controls,
    \item synthetic dataset generation from a fixed cell-type reference dictionary,
    \item existing spot encoder + cross-attention + output head wiring,
    \item a minimal experiment script that trains and writes metrics to JSON.
\end{enumerate}

\section{Mathematical Role of $Q$, $K$, and $V$}
For one batch of spots, let:
\begin{itemize}
    \item spot features be $X \in \mathbb{R}^{B \times G}$,
    \item reference dictionary be $R \in \mathbb{R}^{B \times C \times D_r}$,
    \item model width be $d$.
\end{itemize}

The encoders produce:
\begin{align}
Q &= f_{\text{spot}}(X) \in \mathbb{R}^{B \times 1 \times d},\\
K &= f_{\text{ref}}(R) \in \mathbb{R}^{B \times C \times d},\\
V &= f_{\text{ref}}(R) \in \mathbb{R}^{B \times C \times d}.
\end{align}

Cross-attention computes a context vector for each spot query:
\begin{equation}
\text{Attn}(Q,K,V) = \text{softmax}\left(\frac{QK^\top}{\sqrt{d_h}}\right)V,
\end{equation}
with multi-head splitting over $d_h=d/H$ and then concatenation/projection by
\texttt{MultiheadAttention}. Intuitively, the spot query asks \emph{which cell-type dictionary
entries are most relevant}, and values aggregate those entries into a spot-conditioned latent.

The output head maps the attended representation $Z \in \mathbb{R}^{B \times 1 \times d}$
into cell-type probabilities:
\begin{equation}
\hat{y}=\text{softmax}(WZ+b) \in \mathbb{R}^{B \times C}.
\end{equation}
Training uses KL divergence to match target compositions sampled from a Dirichlet prior.

\section{Synthetic Data Assumptions and Tensor Shapes}
\begin{itemize}
    \item A global reference dictionary $R_0 \in \mathbb{R}^{C \times D_r}$ is sampled once per dataset split.
    \item Spot compositions $y_i \sim \text{Dirichlet}(\alpha\mathbf{1})$.
    \item Clean spot features are generated by linear mixing: $x_i = y_i R_0$ (with optional random projection when $D_r \neq G$), plus Gaussian noise.
\end{itemize}

Default key shapes:
\begin{center}
\begin{tabular}{ll}
\toprule
Tensor & Shape \\
\midrule
spot\_features & $[B, G]$ \\
reference\_embeddings & $[B, C, D_r]$ \\
target\_composition & $[B, C]$ \\
attention weights (averaged heads) & $[B, 1, C]$ \\
\bottomrule
\end{tabular}
\end{center}

\section{Next Steps}
\begin{enumerate}
    \item Replace synthetic dictionary generation with a real cell-type reference matrix from scRNA-seq.
    \item Add optional spatial neighborhood covariates into the spot encoder.
    \item Add calibration/ablation metrics (e.g., JS divergence, per-cell-type AUROC where labels exist).
    \item Benchmark against classical deconvolution baselines on matched datasets.
\end{enumerate}

\end{document}
